\documentclass[14pt,a4paper]{report}
\usepackage[utf8]{inputenc}
\usepackage[french]{babel}
\usepackage{graphicx}
\usepackage{hyperref}
\usepackage{listings}
\usepackage{color}
\usepackage{enumitem}
\usepackage{fancyhdr}
\usepackage{geometry}

% Configuration de la géométrie de la page
\geometry{
    a4paper,
    top=2.5cm,
    bottom=2.5cm,
    left=2.5cm,
    right=2.5cm
}

% Configuration des en-têtes et pieds de page
\pagestyle{fancy}
\fancyhf{}
\rhead{Mercato}
\lhead{Document de Maintenance}
\rfoot{Page \thepage}

% Définition des couleurs pour le code
\definecolor{codegreen}{rgb}{0,0.6,0}
\definecolor{codegray}{rgb}{0.5,0.5,0.5}
\definecolor{codepurple}{rgb}{0.58,0,0.82}
\definecolor{backcolour}{rgb}{0.95,0.95,0.92}

% Configuration des listings pour le code
\lstdefinestyle{mystyle}{
    backgroundcolor=\color{backcolour},
    commentstyle=\color{codegreen},
    keywordstyle=\color{magenta},
    numberstyle=\tiny\color{codegray},
    stringstyle=\color{codepurple},
    basicstyle=\ttfamily\footnotesize,
    breakatwhitespace=false,
    breaklines=true,
    captionpos=b,
    keepspaces=true,
    numbers=left,
    numbersep=5pt,
    showspaces=false,
    showstringspaces=false,
    showtabs=false,
    tabsize=2
}
\lstset{style=mystyle}

\begin{document}

\begin{titlepage}
    \begin{center}
        {\large REPUBLIQUE DU CAMEROUN \hfill REPUBLIC OF CAMEROON}\\
        {\small Paix-Travail-Patrie \hfill Peace-Work-Fatherland}\\[0.5cm]
        {\large UNIVERSITE DE YAOUNDE 1 \hfill UNIVERSITY OF YAOUNDE 1}\\
        {\large FACULTE DES SCIENCES \hfill FACULTY OF SCIENCE}\\
        {\large Département d'Informatique \hfill Department of computer science}\\
        {\small B.P.812 Yaoundé \hfill PO.BOX 812 Yaoundé}\\
        {\small Tél /Fax : (+237) 22 23 44 96 \hfill Tel /Fax : (+237) 22 23 44 96}\\[2cm]

        \begin{center}
            {\large \textbf{Filière} : Informatique \hspace{2cm} \textbf{Niveau} : M1}\\
            {\large \textbf{Unité d'enseignement (UE)} : INF 4027 : Génie Logiciel}\\[2cm]

            \framebox{
                \begin{minipage}{0.8\textwidth}
                    \begin{center}
                        \huge \textbf{MERCATO}\\[0.5cm]
                        \Large Document de Maintenance
                    \end{center}
                \end{minipage}
            }\\[2cm]

            \textbf{Rédigé par} :\\[0.5cm]
            \begin{tabular}{|l|l|}
                \hline
                \textbf{Noms \& Prénoms} & \textbf{Matricule} \\
                \hline
                MAHAMAT SALEH MAHAMAT & 21T2423 \\
                \hline
                SOKOUDJOU CHENDOU CHRISTIAN MANUEL & 21T2396 \\
                \hline
                ZOGO ABOUMA ZOZIME ACHAIRE & 18N2824 \\
                \hline
                FEUJO TSAGMO SIMPLICE JORDA & 18T2897 \\
                \hline
            \end{tabular}\\[2cm]

            \textbf{Supervisé par} :\\
            Pr. Roger ATSA ETOUNDI\\[1cm]

            \textbf{Année universitaire} : 2024 - 2025
        \end{center}
    \end{center}
\end{titlepage}

\tableofcontents
\newpage

\chapter{Système de Paiement}

\section{Surveillance des Transactions}
\subsection{Monitoring en Temps Réel}
Le système de monitoring en temps réel est basé sur les composants suivants :
\begin{itemize}
    \item Dashboard de surveillance des transactions
    \item Système de métriques avec Prometheus
    \item Visualisation Grafana pour les KPIs
    \item Seuils d'alerte configurés :
        \begin{itemize}
            \item Taux d'erreur > 5\%
            \item Temps de réponse > 2s
            \item Volume de transactions anormal
        \end{itemize}
\end{itemize}

\subsection{Système d'Alertes}
Configuration des alertes via :
\begin{itemize}
    \item Intégration Slack pour les notifications immédiates
    \item Emails pour les rapports détaillés
    \item SMS pour les incidents critiques
\end{itemize}

\subsection{Gestion des Logs}
Procédures de gestion des logs :
\begin{itemize}
    \item Rotation quotidienne des logs
    \item Rétention de 90 jours
    \item Agrégation via ELK Stack
    \item Analyse automatisée des patterns d'erreur
\end{itemize}

\section{Maintenance des Intégrations}
\subsection{Système de Carte Bancaire}
\subsubsection{Certification SSL}
\begin{itemize}
    \item Vérification mensuelle des dates d'expiration
    \item Procédure de renouvellement :
        \begin{enumerate}
            \item Génération de la CSR
            \item Validation par l'autorité de certification
            \item Installation et test du nouveau certificat
            \item Période de chevauchement de 7 jours
        \end{enumerate}
\end{itemize}

\subsubsection{Tests de Transaction}
Programme de tests automatisés :
\begin{itemize}
    \item Tests quotidiens des transactions de base
    \item Tests hebdomadaires des cas limites
    \item Simulation mensuelle de charge importante
\end{itemize}

\subsection{Mobile Money}
\subsubsection{MTN Money}
Maintenance spécifique :
\begin{itemize}
    \item Surveillance du taux de succès des transactions
    \item Gestion des tokens d'authentification
    \item Monitoring des limites de transaction
\end{itemize}

\subsubsection{Orange Money}
Points de contrôle :
\begin{itemize}
    \item Vérification des webhooks
    \item Gestion des timeouts (configuré à 30s)
    \item Surveillance des soldes des comptes marchands
\end{itemize}

\chapter{Maintenance de la Base de Données}

\section{Optimisation}
\subsection{Nettoyage des Données}
Procédures automatisées :
\begin{itemize}
    \item Archivage des transactions > 1 an
    \item Suppression des sessions expirées
    \item Compression des logs anciens
\end{itemize}

\subsection{Performance}
Optimisations régulières :
\begin{itemize}
    \item VACUUM ANALYZE hebdomadaire
    \item Mise à jour des statistiques
    \item Révision des index
    \item Monitoring des requêtes lentes
\end{itemize}

\section{Stratégie de Sauvegarde}
\subsection{Backup Incrémental}
Configuration :
\begin{itemize}
    \item Fréquence : toutes les 6 heures
    \item Rétention : 7 jours
    \item Compression : niveau 6
    \item Vérification automatique de l'intégrité
\end{itemize}

\subsection{Backup Complet}
Procédure hebdomadaire :
\begin{itemize}
    \item Exécution : dimanche 02:00
    \item Rétention : 3 mois
    \item Stockage externe chiffré
    \item Test de restauration mensuel
\end{itemize}

\chapter{Maintenance Applicative}

\section{Frontend}
\subsection{Optimisation des Performances}
Actions régulières :
\begin{itemize}
    \item Minification des assets
    \item Compression des images
    \item Mise en cache des ressources statiques
    \item Lazy loading des composants
\end{itemize}

\subsection{Compatibilité}
Tests systématiques sur :
\begin{itemize}
    \item Chrome, Firefox, Safari, Edge
    \item Versions mobiles
    \item Différentes résolutions d'écran
\end{itemize}

\section{Backend}
\subsection{Gestion des Dépendances}
Procédures de mise à jour :
\begin{itemize}
    \item Analyse hebdomadaire des vulnérabilités
    \item Tests de régression automatisés
    \item Documentation des changements
\end{itemize}

\subsection{Cache}
Configuration et maintenance :
\begin{itemize}
    \item Redis pour le cache applicatif
    \item Memcached pour les sessions
    \item Politique de purge automatique
\end{itemize}

\chapter{Procédures d'Urgence}

\section{Gestion des Incidents}
\subsection{Niveaux de Criticité}
\begin{itemize}
    \item Niveau 1 : Impact mineur
    \item Niveau 2 : Service dégradé
    \item Niveau 3 : Service interrompu
    \item Niveau 4 : Incident critique
\end{itemize}

\subsection{Procédure de Rollback}
Étapes à suivre :
\begin{enumerate}
    \item Évaluation de l'impact
    \item Sauvegarde de l'état actuel
    \item Exécution du rollback
    \item Vérification post-rollback
    \item Communication aux parties prenantes
\end{enumerate}

\section{Documentation des Incidents}
Informations à consigner :
\begin{itemize}
    \item Date et heure de l'incident
    \item Nature du problème
    \item Actions entreprises
    \item Impact sur le service
    \item Mesures préventives identifiées
\end{itemize}

\chapter{Coûts de Maintenance}

\section{Coûts Récurrents}
\begin{table}[h]
    \centering
    \begin{tabular}{|l|r|p{6cm}|}
        \hline
        \textbf{Poste} & \textbf{Coût Mensuel (XAF)} & \textbf{Description} \\
        \hline
        Support Technique & 100.000 & Assistance technique niveau 1 et 2 \\
        \hline
        Maintenance Préventive & 75.000 & Vérifications et mises à jour régulières \\
        \hline
        Maintenance Corrective & 50.000 & Résolution des bugs et incidents \\
        \hline
        Sauvegardes & 25.000 & Gestion des backups et restaurations \\
        \hline
        \textbf{Total} & \textbf{250.000} & \textit{Coût mensuel} \\
        \hline
    \end{tabular}
    \caption{Coûts Mensuels de Maintenance}
\end{table}

\section{Coûts Exceptionnels}
\begin{table}[h]
    \centering
    \begin{tabular}{|l|r|p{6cm}|}
        \hline
        \textbf{Intervention} & \textbf{Coût (XAF)} & \textbf{Fréquence} \\
        \hline
        Mise à jour majeure & 300.000 & Semestrielle \\
        \hline
        Audit sécurité & 200.000 & Annuel \\
        \hline
        Formation continue & 150.000 & Trimestrielle \\
        \hline
    \end{tabular}
    \caption{Coûts des Interventions Exceptionnelles}
\end{table}

\chapter{Environnements}

\section{Environnement de Préproduction}
\subsection{Configuration Matérielle}
\begin{itemize}
    \item Serveur dédié : Dell PowerEdge T150 (ou équivalent)
    \item Processeur : Intel Xeon E-2324G
    \item RAM : 16 GB
    \item Stockage : 2 x 1TB en RAID 1
    \item Système d'exploitation : Ubuntu Server 22.04 LTS
\end{itemize}

\subsection{Configuration Logicielle}
\begin{itemize}
    \item Base de données : MySQL 8.0
    \item Serveur Web : Nginx
    \item PHP 8.1
    \item Node.js 18 LTS
    \item Git pour le versioning
\end{itemize}

\section{Environnement de Formation}
\subsection{Configuration}
\begin{itemize}
    \item Poste de formation dédié
    \item Base de données de test
    \item Jeu de données de formation
    \item Documentation utilisateur
    \item Guides pratiques
\end{itemize}

\subsection{Ressources}
\begin{itemize}
    \item Salle de formation équipée
    \item Support de formation
    \item Exercices pratiques
    \item Scénarios de test
\end{itemize}

\section{Mise en Production}
\subsection{Procédure de Déploiement}
\begin{enumerate}
    \item \textbf{Préparation}
        \begin{itemize}
            \item Vérification des prérequis
            \item Sauvegarde complète
            \item Information des utilisateurs
        \end{itemize}
    \item \textbf{Installation}
        \begin{itemize}
            \item Déploiement des fichiers
            \item Configuration du serveur
            \item Migration des données
        \end{itemize}
    \item \textbf{Tests}
        \begin{itemize}
            \item Tests fonctionnels
            \item Tests de performance
            \item Validation utilisateur
        \end{itemize}
    \item \textbf{Mise en service}
        \begin{itemize}
            \item Activation du système
            \item Formation des utilisateurs
            \item Période de surveillance
        \end{itemize}
\end{enumerate}

\subsection{Plan de Bascule}
\begin{itemize}
    \item \textbf{Timing} : Week-end (48h)
    \item \textbf{Équipe mobilisée} :
        \begin{itemize}
            \item 1 Chef de projet
            \item 2 Techniciens
            \item 1 Formateur
        \end{itemize}
    \item \textbf{Plan de secours} :
        \begin{itemize}
            \item Procédure de rollback
            \item Restauration des sauvegardes
            \item Support d'urgence
        \end{itemize}
\end{itemize}

\chapter{Annexes}

\section{Contacts d'Urgence}
\begin{itemize}
    \item Support Niveau 1 : +XXX XXX XXX
    \item Support Niveau 2 : +XXX XXX XXX
    \item Administrateur Système : +XXX XXX XXX
    \item Responsable Sécurité : +XXX XXX XXX
\end{itemize}

\section{Outils de Monitoring}
\begin{itemize}
    \item Dashboard principal : \url{http://monitoring.mercato.local}
    \item Logs : \url{http://logs.mercato.local}
    \item Métriques : \url{http://metrics.mercato.local}
\end{itemize}

\end{document}
